\chapter{Introducció i Context}
\label{ch:introduccio}

En aquest primer capítol es defineix el marc en el qual es desenvolupa el projecte, s'analitza la problemàtica de la plataforma web \textit{legacy} de l'empresa Lanaccess i es justifica la necessitat de substituir-la per una solució moderna. Així mateix, es detallen els objectius, l'abast i els actors implicats en el cicle de vida del software.

\section{Context del projecte}
Aquest Treball de Final de Grau (TFG) consisteix en el disseny i desenvolupament de \textbf{Web Core View}, una aplicació web moderna destinada a la gestió centralitzada dels videogravadors (\textit{Network Video Recorder}, d'ara endavant \textbf{NVR}) de la família \textbf{NXT} de l'empresa Lanaccess Telecom, SA. El projecte es realitza en col·laboració amb aquesta empresa tecnològica, especialitzada en el desenvolupament de solucions professionals de videovigilància i seguretat electrònica.

El treball s'emmarca dins de la \textbf{Modalitat B} del TFG i de l'especialitat d'Enginyeria de Software, posant èmfasi en arquitectures \textit{frontend} i \textit{backend} modernes, seguretat en la comunicació i integració de protocols. L'objectiu és modernitzar la plataforma tecnològica de l'empresa, substituint l'aplicació web \textit{legacy} construïda als anys 90 amb tecnologies avui obsoletes (\textbf{ActiveX}) per una interfície de control i administració \textit{embedded} que s'executa directament al propi videogravador NXT, basada en l'ecosistema Angular i protocols estàndard oberts (WebRTC, HLS, WSS).

\section{Identificació del problema i motivació}
La plataforma \textit{legacy} de Lanaccess és una aplicació web construïda als anys 90 que depenia en la seva totalitat dels controls \textbf{ActiveX} de Microsoft per a la visualització de vídeo i la interacció amb el maquinari. Amb la retirada oficial del suport a Internet Explorer 11 per part de Microsoft el juny del 2022 \cite{microsoft_ie11_eol}, l'últim navegador capaç d'executar controls ActiveX va quedar obsolet, fent la plataforma pràcticament inoperativa per a qualsevol instal·lació moderna. Aquesta obsolescència es concreta en tres problemes fonamentals:

\begin{itemize}
  \item \textbf{Incompatibilitat amb navegadors moderns:} ActiveX és una tecnologia exclusiva d'Internet Explorer. Els navegadors actuals (Chrome, Firefox, Edge, Safari) han eliminat completament el suport per a \textit{plugins} propietaris externs per raons de seguretat, de manera que la plataforma \textit{legacy} és avui inaccessible des de qualsevol navegador d'ús estàndard.
  \item \textbf{Rigidesa davant hardware heterogeni:} La interfície estàtica de la plataforma \textit{legacy} no pot adaptar-se a la disparitat de paràmetres i \textit{firmwares} dels videogravadors NXT (adreces IP, ports, lògiques d'alarma, còdecs suportats). Qualsevol canvi en el maquinari suportat requereix modificar i redistribuir manualment el codi de l'aplicació.
  \item \textbf{Absència de capacitats en temps real:} La interfície, basada en pàgines estàtiques dels anys 90 (recàrregues completes, sense comunicació asíncrona), no ofereix cap de les funcionalitats de temps real que requereixen els operadors de seguretat actuals: visualització de vídeo en directe, alertes instantànies ni configuració dinàmica.
\end{itemize}

Paral·lelament, Lanaccess ha identificat una \textbf{necessitat de mercat addicional}: oferir una solució de gestió autònoma orientada a \textbf{empreses petites} amb un únic videogravador NXT. El model anterior era una aplicació web centralitzada, allotjada en un servidor intern de la instal·lació, que gestionava simultàniament tot el parc de dispositius d'una gran infraestructura. Aquesta arquitectura resulta sobredimensionada i innecessàriament complexa per a clients amb una sola unitat de gravació. Així, \textit{Web Core View} neix amb una doble finalitat: modernitzar la plataforma tecnològica obsoleta i, alhora, proporcionar una solució \textit{embedded} lleugera que s'executi directament en el propi videogravador NXT, permetent la gestió autònoma de l'equip des de qualsevol navegador sense requerir infraestructura centralitzada addicional.

\section{Estat de l'art i alternatives existents}

\subsection{Estat de l'art}
El mercat de la videovigilància IP és un dels sectors tecnològics de major creixement, en el qual fabricants com Hikvision i Dahua han assolit una posició de lideratge gràcies a dispositius d'alt rendiment i cost competitiu \cite{ihs_video_surveillance}. Les seves solucions incorporen interfícies web integrades als videogravadors, però sovint requereixen la instal·lació de \textit{plugins} propietaris per a la visualització de vídeo, fet que limita la seva compatibilitat a navegadors específics i planteja riscos de seguretat corporativa.

Aquesta limitació reflecteix una tendència de canvi estructural del sector cap al \textbf{vídeo \textit{plugin-free}}. L'aparició d'estàndards oberts com \textbf{WebRTC} \cite{rfc8825} i \textbf{HLS} \cite{hls_apple} ha suposat un punt d'inflexió: permeten la transmissió de fluxos multimèdia de baixa latència i la reproducció d'enregistraments de forma nativa al navegador, sense extensions externes. Paral·lelament, la gran majoria d'organitzacions bloquegen per política interna la instal·lació de \textit{plugins} a les estacions de treball dels operadors per evitar vectors d'atac, fet que converteix el \textit{plugin-free} en un requisit de compliment de seguretat corporativa per a grans infraestructures.

D'altra banda, les plataformes VMS (\textit{Video Management System}) de gamma alta, com Milestone XProtect \cite{milestone_xprotect} o Genetec Security Center \cite{genetec_security_center}, ofereixen solucions modulars i molt potents, però requereixen una infraestructura de servidors dedicada, costos de llicenciament per canal elevats i personal TI especialitzat. Estan orientades a grans instal·lacions corporatives i no contemplen un model de desplegament autònom en un únic NVR sense servidors addicionals.

\subsection{Taula comparativa}
La Taula~\ref{tab:comparativa_mercat} compara les dues categories de solucions de mercat predominants, la plataforma \textit{legacy} i la proposta de Web Core View, diferenciant entre fabricants de maquinari (Hikvision/Dahua) i plataformes VMS d'empresa (Milestone/Genetec).

\begin{table}[h!]
  \centering
  \caption{Comparativa de solucions de videovigilància web.}
  \label{tab:comparativa_mercat}
  \resizebox{\textwidth}{!}{
    \begin{tabular}{@{}lllll@{}}
      \toprule
      \textbf{Característica} & \textbf{Web Legacy (ActiveX)} & \textbf{Hikvision / Dahua} & \textbf{VMS (Milestone / Genetec)} & \textbf{Web Core View} \\ \midrule
      Accés web          & IE / ActiveX (obsolet)  & Parcial (\textit{plugins}) & App dedicada / client ric     & \textbf{Navegador modern}              \\
      Protocol de vídeo  & RTSP via ActiveX        & RTSP / H.264               & RTSP / RTMP                   & \textbf{WebRTC / HLS}                  \\
      Configuració       & Estàtica                & Fixa per model             & Propietària                   & \textbf{Dinàmica (\textit{Schema-driven})} \\
      Cost / llicències  & N/A (intern)            & Inclòs en hardware         & Alt (per canal/càmera)        & \textbf{Inclòs al NXT}                 \\
      Desplegament       & Servidor centralitzat   & Servidor centralitzat      & Servidors dedicats            & \textbf{Autònom (\textit{embedded})}   \\ \bottomrule
    \end{tabular}
  }
  \par\smallskip
  \small Font: Elaboració pròpia a partir de la documentació pública dels fabricants \cite{milestone_xprotect, genetec_security_center}.
\end{table}

\section{Anàlisi d'alternatives i justificació de la solució}
Per garantir l'èxit i la sostenibilitat del projecte, s'han avaluat diferents camins tecnològics per als pilars fonamentals de programari, vídeo, concurrència i em\-ma\-gat\-ze\-mat\-ge del sistema:

\subsection{Alternativa de Framework Frontend: Angular 19 vs. React / Vue}
\begin{itemize}
  \item \textbf{React / Vue:} Són llibreries extremadament populars i flexibles, usades en entorns tan exigents com la banca o l'administració pública. En ser arquitectures no opinades, requereixen integrar múltiples llibreries de tercers (ruteig, formularis, \textit{state management}) per cobrir funcionalitats que en un \textit{framework} empresarial ja venen integrades, cosa que implica un cost de manteniment i auditoria de dependències superior en un projecte de vida llarga.
  \item \textbf{Angular 19:} És un \textit{framework} empresarial complet: proporciona ruteig, formularis reactius, injecció de dependències i un sistema de mòduls propi. Escrit nativament en TypeScript, incorpora eines de \textit{testing} integrades (Jasmine/Karma) i una estructura opinada que facilita la consistència del codi en equips.
  \item \textbf{Justificació:} S'ha escollit \textbf{Angular 19} per la seva estructura opinada que redueix decisions arquitecturals arbitràries, el tipatge estàtic de TypeScript, el motor de formularis reactius nadiu, les eines de \textit{testing} integrades i, principalment, la introducció de les \textbf{Signals} \cite{angular_signals} com a mecanisme natiu i eficient de gestió de l'estat reactiu, que permet prescindir de \textit{Zone.js} i optimitza el rendiment en dispositius \textit{embedded}.
\end{itemize}

\subsection{Alternativa de vídeo: Arquitectura Híbrida WebRTC i HLS vs. RTSP / RTMP}
\begin{itemize}
  \item \textbf{RTSP (Real-Time Streaming Protocol):} És el protocol de facto utilitzat internament per les càmeres IP i sistemes VMS clàssics. No obstant això, és completament incompatible de forma nativa amb els navegadors web actuals, obligant a utilitzar transcodificadors molt costosos o \textit{plugins} externs propietaris en desús.
  \item \textbf{RTMP (Real-Time Messaging Protocol):} Protocol clàssic molt lligat a la tecnologia ja obsoleta d'Adobe Flash, actualment completament desapareguda de l'ecosistema web estàndard.
  \item \textbf{WebRTC i HLS:} Són els dos únics estàndards web moderns reconeguts pel W3C que permeten la reproducció de vídeo nativa (HTML5) de forma segura i sense cap mena d'extensió o control extern.
  \item \textbf{Justificació:} S'ha optat per una \textbf{arquitectura híbrida complementària} que explota el millor de cada protocol en lloc de triar-ne només un:
        \begin{enumerate}
          \item \textbf{WebRTC:} S'utilitza principalment per al \textbf{vídeo en directe}, on la latència ultra-baixa ($<300$~ms) i el transport sobre UDP sota xifratge DTLS són indispensables per vigilar o respondre a amenaces immediates.
          \item \textbf{HLS (HTTP Live Streaming) \cite{hls_apple}:} S'utilitza per a la \textbf{reproducció de gravacions històriques}, on la ultra-baixa latència no és prioritària però es requereix un control molt precís de descàrrega per blocs (segments), commutació dinàmica de qualitat i la capacitat de fer desplaçaments lliures (\textit{seek}) sobre l'històric mitjançant TCP. També actua com a \textit{fallback} per al vídeo en directe si els tallafocs bloquegen el trànsit WebRTC.
        \end{enumerate}
\end{itemize}

\subsection{Alternativa de Base de Dades per a Alarmes: SQLite3 vs. PostgreSQL / MongoDB}
\begin{itemize}
  \item \textbf{PostgreSQL / MySQL:} Són bases de dades relacionals molt potents, però funcionen com a serveis de sistema independents. Requereixen un consum elevat i constant de memòria RAM i CPU per mantenir el dimoni actiu, sent una sobrecàrrega inacceptable en dispositius embedded limitats.
  \item \textbf{MongoDB:} Base de dades documental NoSQL ideal per a grans volums de dades desestructurades, però amb una pila tecnològica pesada i requisits de recursos inassequibles per al maquinari de gravació local.
  \item \textbf{SQLite3:} És un motor de base de dades relacional \textit{serverless}, transaccional (ACID) i auto-tingut que s'executa en el mateix procés del backend (in-process). Guarda tota la informació en un únic fitxer físic al disc.
  \item \textbf{Justificació:} S'ha seleccionat \textbf{SQLite3} com el motor de base de dades per al registre d'alarmes i logs del dispositiu, atès que ofereix un rendiment excel·lent per a lectures i escriptures ràpides amb un consum residual de memòria ($<15$~MB), essent ideal per als gravadors NXT de Lanaccess amb recursos de memòria i CPU limitats.
\end{itemize}

\subsection{Alternativa de Còdec de Vídeo: H.264 i H.265 coexistents}
\begin{itemize}
  \item \textbf{H.264 (AVC):} Còdec universal amb descodificació hardware garantida en el 100\% dels navegadors i dispositius. Consum d'ample de banda i d'espai de disc superior per a la mateixa qualitat d'imatge.
  \item \textbf{H.265 (HEVC):} Redueix fins a la meitat l'espai d'emmagatzematge i l'ample de banda per a la mateixa qualitat, però el seu suport web depèn de la disponibilitat de descodificació hardware al terminal client, degut a restriccions de patents.
  \item \textbf{Justificació:} S'ha optat per una \textbf{estratègia dinàmica de coexistència}: el sistema prioritza H.265 quan la càmera i el terminal client ho suporten (màxim estalvi de disc), i commuta automàticament a H.264 en cas contrari (portabilitat universal). Els detalls de la detecció de capacitats i la lògica de commutació s'expliquen al Capítol~\ref{ch:implementacio}.
\end{itemize}

\subsection{Alternativa de configuració: Estàtica vs. Schema-driven}
\begin{itemize}
  \item \textbf{Formularis estàtics:} Programar una pantalla per a cada menú de configuració és més ràpid a curt termini però genera un deute tècnic insostenible: cada canvi en el \textit{firmware} del videogravador NXT —un nou paràmetre, un rang vàlid diferent, un camp eliminat— obligaria a actualitzar i redistribuir el \textit{frontend}. En una plataforma amb múltiples models de hardware i \textit{firmwares} en evolució continuada, això és inviable.
  \item \textbf{Motor Schema-driven:} El \textit{backend} envia al \textit{frontend} un diccionari pla de claus en notació de punts (p. ex., \texttt{camera.\{id\}.stream.codec}) on cada valor porta metadades amb prefix \textit{underscore}: \texttt{\_type}, \texttt{\_allowedValues}, \texttt{\_min}, \texttt{\_max}, \texttt{\_default} i \texttt{\_description}. El \textit{frontend} actua com un intèrpret universal que renderitza automàticament el control adequat —selector, camp numèric, \textit{toggle}, etc.— sense conèixer el model concret de videogravador ni la versió de \textit{firmware}.
  \item \textbf{Justificació:} S'ha optat per la solució \textbf{Schema-driven} perquè és la característica diferenciadora del projecte: permet que \textbf{Web Core View} sigui compatible amb qualsevol model futur de la família NXT de Lanaccess sense modificar ni una línia de codi del \textit{frontend}. El format real del \textit{schema} s'explica en detall a l'Annex~A.1 de la memòria.
\end{itemize}

\subsection{Alternativa de gestió d'estat: NgRx vs. Angular Signals}
\begin{itemize}
  \item \textbf{NgRx:} Implementació del patró Redux per a l'ecosistema Angular, robusta per a aplicacions gegantines amb estats complexos compartits entre mòduls. Implica un \textit{boilerplate} elevat (accions, \textit{reducers}, efectes, selectors) i un cost d'aprenentatge notable.
  \item \textbf{Angular Signals:} La nova eina nativa d'Angular 19 per a la gestió de la reactivitat granular \cite{angular_signals}. Permet que els components reaccionin exactament als canvis d'estat que els afecten, sense necessitat de \textit{Zone.js}.
  \item \textbf{Justificació:} S'ha seleccionat \textbf{Signals} per la seva lleugeresa i eficiència en el renderitzat. En un entorn \textit{embedded} on la CPU s'ha de prioritzar per al vídeo, la reactivitat nativa de Signals permet prescindir de les comprovacions constants de \textit{Zone.js}, optimitzant el rendiment global del sistema.
\end{itemize}

\subsection{Alternativa de mecanisme de sessió: JWT vs. Cookies HttpOnly/Secure}
\begin{itemize}
  \item \textbf{JWT (\textit{JSON Web Token}):} Mecanisme d'autenticació sense estat (\textit{stateless}) \cite{rfc7519} que codifica les credencials al propi \textit{token} signat. La gestió des del \textit{frontend} implica emmagatzemar el \textit{token} a \textit{LocalStorage} o \textit{SessionStorage}, fet que els exposa a atacs XSS (\textit{Cross-Site Scripting}) si hi ha una injecció de codi maliciós.
  \item \textbf{Cookies \texttt{HttpOnly}/\texttt{Secure}:} La \textit{cookie} de sessió és gestionada exclusivament per NGINX: s'estableix al \textit{backend}, s'envia al navegador amb els atributs \texttt{HttpOnly} i \texttt{Secure}, i el JavaScript de l'aplicació no en té accés en cap moment. L'estat de sessió es manté en un Signal privat en memòria al \textit{frontend} (zero-persistència: un F5 equival a tancar sessió).
  \item \textbf{Justificació:} S'ha optat per les \textbf{cookies \texttt{HttpOnly}/\texttt{Secure}} gestionades per NGINX perquè eliminen el vector d'atac XSS sobre les credencials: cap \textit{script} del navegador no pot llegir ni exfiltrar la \textit{cookie}. Aquesta decisió s'alinea amb les recomanacions de seguretat de l'OWASP per a aplicacions en entorns de control d'accés físic.
\end{itemize}

\subsection{Alternativa de notificacions en temps real: WebSockets vs. SSE vs. Long Polling}
\begin{itemize}
  \item \textbf{Long Polling:} El client fa peticions HTTP repetides fins que el servidor disposa de dades noves. Solució senzilla però ineficient: genera múltiples connexions curtes i latències variables, inadequada per a notificacions d'alarmes crítiques.
  \item \textbf{SSE (\textit{Server-Sent Events}):} Connexió HTTP unidireccional (servidor $\rightarrow$ client), lleugera i nativa al navegador. Adequada per a \textit{feeds} de dades en un sol sentit, però no permet enviar comandes del client al servidor per la mateixa connexió.
  \item \textbf{WebSockets (WSS):} Connexió TCP persistent i bidireccional \cite{rfc6455} que permet tant rebre notificacions \textit{push} del servidor com enviar comandes del client (p. ex., marcar una alarma com a llegida) amb latència mínima.
  \item \textbf{Justificació:} S'ha escollit \textbf{WebSockets} per la seva naturalesa bidireccional, necessària per als dos canals del sistema: el canal global de notificacions (\texttt{/ws/la-core-no\-ti\-fier/}) i el canal d'alarmes (\texttt{/ws/la-core-alarm/}), que requereix tant rebre alertes com confirmar-ne la lectura des del client.
\end{itemize}

\subsection{Alternativa de proxy invers: NGINX vs. accés directe als microserveis}
\begin{itemize}
  \item \textbf{Accés directe als microserveis:} Exposar cada microservei del \textit{backend} en ports independents obliga el navegador a realitzar peticions \textit{cross-origin}, fent necessaris encapçalaments CORS complexos i impedint l'ús de cookies \texttt{HttpOnly} (que no es transmeten en peticions \textit{cross-origin} per defecte).
  \item \textbf{NGINX com a proxy invers:} Centralitza tots els microserveis sota un únic origen (\textit{same-origin}), enrutant per prefix de ruta. Gestiona el cicle de vida de les \textit{cookies} de sessió de forma transparent i actua com a punt únic de control de seguretat (TLS, capçaleres HTTP de seguretat).
  \item \textbf{Justificació:} S'ha escollit \textbf{NGINX} \cite{nginx_admin} per la seva maduresa i rendiment en sistemes Linux embeguts: permet centralitzar la seguretat, implementar les \textit{cookies} \texttt{HttpOnly}/\texttt{Secure} de forma consistent sense modificar cap microservei del \textit{backend}, i elimina el problema CORS de manera estructural.
\end{itemize}

\section{Solució proposada i valor afegit}
\textbf{Web Core View} es proposa com una solució integral que aprofita les noves funcionalitats de rendiment d'\textbf{Angular 19}. L'elecció d'aquesta versió respon a la voluntat d'utilitzar tecnologies actualitzades i de llarg suport. Les aportacions clau de la solució proposada són:
\begin{itemize}
  \item \textbf{Arquitectura Schema-driven:} La interfície interpreta un esquema JSON enviat pel gravador i genera automàticament els controls de configuració, assegurant compatibilitat amb qualsevol versió de hardware.
  \item \textbf{Gestió d'estat amb Signals:} L'ús d'\textbf{Angular Signals} representa un canvi de paradigma en la reactivitat. Aquesta funcionalitat permet una gestió granular dels canvis, optimitzant el rendiment del \textit{frontend} i garantint una fluïdesa absoluta fins i tot amb múltiples fluxos de vídeo simultanis, evitant sobrecàrregues de CPU innecessàries al terminal del client.
  \item \textbf{Control de concurrència:} Implementació d'un servei de bloqueig en temps real per evitar conflictes d'edició entre administradors.
  \item \textbf{Arquitectura d'esdeveniments:} Ús de \textit{WebSockets} (WSS) per rebre notificacions d'alarmes, notificacions del sistema i errors de forma immediata.
  \item \textbf{Desplegament autònom (single-NVR):} L'aplicació s'executa directament dins del propi videogravador NXT com a servidor web \textit{embedded}. Això permet que empreses petites amb un únic dispositiu puguin gestionar-lo de forma completa i autònoma des de qualsevol navegador, sense necessitat d'infraestructura centralitzada addicional ni de programari extern. Aquesta capacitat obre un nou segment de mercat per a Lanaccess que la plataforma \textit{legacy} anterior no podia cobrir.
\end{itemize}

\section{Objectius del projecte}
\subsection{Objectiu general}
Dissenyar i desenvolupar una consola de gestió web modular, reactiva i d'alt rendiment que centralitzi la gestió tècnica i operativa dels NVRs NXT de Lanaccess, substituint la plataforma web \textit{legacy} dels anys 90 i les seves tecnologies obsoletes (ActiveX) per una arquitectura i interfície completament modernes.

\subsection{Objectius específics}
\begin{itemize}
  \item Desenvolupar un motor de configuració dinàmic basat en esquemes metadada.
  \item Integrar un mòdul de visualització híbrid compatible amb WebRTC i HLS (H264/H265).
  \item Implementar la gestió d'alarmes en temps real mitjançant comunicació bidireccional.
  \item Facilitar el manteniment remot (\textit{firmware}, llicències i logs) via navegador.
  \item Dissenyar l'aplicació com una solució autònoma \textit{embedded} capaç de funcionar directament en un únic NVR NXT (\textit{single-NVR}), sense dependència de servidors centralitzats, per cobrir el segment d'empreses petites.
\end{itemize}

\section{Abast i Obstacles}
\subsection{Abast funcional}
L'aplicació ha d'oferir un cicle complet de gestió: des de l'autenticació segura de l'usuari fins a la monitorització de càmeres en directe, la cerca de gravacions històriques i la configuració avançada de xarxa i paràmetres de sistema. El disseny ha de contemplar el desplegament autònom en un únic videogravador NXT (mode \textit{single-NVR}) com a cas d'ús principal, garantint que tota la funcionalitat sigui operable sense infraestructura externa.

\subsection{Obstacles detectats}
S'han identificat reptes tècnics que afecten diferents actors:
\begin{itemize}
  \item \textbf{Complexitat del vídeo natiu:} El repte de visualitzar H.265 sense \textit{plugins}. Afecta al \textbf{Projectista} (recerca de llibreries) i a l'\textbf{Usuari final} (potencial latència).
  \item \textbf{Sincronització en temps real:} Requereix una gestió d'estat robusta perquè l'\textbf{Operador} vegi alarmes sense retard.
  \item \textbf{Seguretat i NGINX:} Configuració del \textbf{Proxy Invers} per centralitzar els microserveis sota un mateix origen i garantir polítiques de \textit{Same-Origin}.
\end{itemize}

\section{Actors implicats (Stakeholders)}
\begin{enumerate}
  \item \textbf{Usuari final (Operador):} Monitoritza càmeres i gestiona incidents.
  \item \textbf{Administrador tècnic:} Configura el sistema i realitza manteniment.
  \item \textbf{Lanaccess Telecom:} Proporciona el \textit{hardware} i coneixement de negoci.
  \item \textbf{Projectista (Estudiant):} Responsable únic de l'anàlisi i la implementació.
  \item \textbf{Equip docent (Director i Ponent):} Supervisen el rigor acadèmic i tècnic.
\end{enumerate}